\documentclass[aspectratio=169,11pt]{beamer}
\usetheme{Madrid}
\usecolortheme{default}
\setbeamertemplate{navigation symbols}{}
\setbeamertemplate{footline}[frame number]
\usepackage{enumitem}

\usepackage{amsmath,amssymb}
\usepackage{booktabs}
\usepackage{array}

\title{Databricks ML Associate}
\subtitle{100 Exam-Style Questions}
\author{Practice Exam Deck}
\date{}

\begin{document}
	
	\begin{frame}
		\titlepage
	\end{frame}
	
	\begin{frame}{How to Use This Deck}
		\begin{itemize}
			\item 100 questions, each on its own slide with four options (A--D).
			\item The \textbf{next slide} gives the correct answer and explains:
			\begin{itemize}
				\item why the correct option is right;
				\item why each wrong option is wrong.
			\end{itemize}
			\item Domains: Databricks ML (Q1--38), ML Workflows (Q39--57), Model Development (Q58--88), Model Deployment (Q89--100).
		\end{itemize}
	\end{frame}
	
	% ---------------------------------------------------------------
	% Q1
	\begin{frame}{Question 1}
		A machine learning engineer created a Feature Table \texttt{new\_table} using Feature Store Client \texttt{fs} and specified a metadata description. Which line of code returns that description?
		\begin{enumerate}[label=\Alph*.]
			\item \texttt{fs.create\_training\_set("new\_table")}
			\item \texttt{fs.get\_table("new\_table").description}
			\item \texttt{fs.get\_table("new\_table").load\_df()}
			\item \texttt{fs.get\_table("new\_table")}
		\end{enumerate}
	\end{frame}
	\begin{frame}{Answer 1}
		\textbf{Correct: B} --- \texttt{get\_table()} returns the table object and \texttt{.description} reads its metadata.\\
		\vspace{0.5em}
		\textbf{Wrong:}
		\begin{itemize}
			\item A creates a training set, not metadata retrieval.
			\item C loads data into a DataFrame.
			\item D returns the object but not the description.
		\end{itemize}
	\end{frame}
	
	% ---------------------------------------------------------------
	\begin{frame}{Question 2}
		What is the primary advantage of creating Feature Store tables at the account level in Unity Catalog rather than at the workspace level?
		\begin{enumerate}[label=\Alph*.]
			\item Faster model training
			\item Cross-workspace governance, discoverability, and access control
			\item Automatic feature engineering
			\item Lower storage cost
		\end{enumerate}
	\end{frame}
	\begin{frame}{Answer 2}
		\textbf{Correct: B} --- account-level tables are governed and discoverable across workspaces.\\
		\vspace{0.5em}
		\textbf{Wrong:}
		\begin{itemize}
			\item A is not a feature of account-level tables.
			\item C is false; feature engineering is manual.
			\item D is not the primary advantage.
		\end{itemize}
	\end{frame}
	
	% ---------------------------------------------------------------
	\begin{frame}{Question 3}
		Which client is used to create and manage Feature Store tables in Unity Catalog?
		\begin{enumerate}[label=\Alph*.]
			\item \texttt{FeatureStoreClient}
			\item \texttt{FeatureEngineeringClient}
			\item \texttt{UnityCatalogClient}
			\item \texttt{MLflowClient}
		\end{enumerate}
	\end{frame}
	\begin{frame}{Answer 3}
		\textbf{Correct: B} --- \texttt{FeatureEngineeringClient} is the Unity Catalog Feature Store client.\\
		\vspace{0.5em}
		\textbf{Wrong:}
		\begin{itemize}
			\item A is the legacy workspace client.
			\item C is not a real client.
			\item D manages MLflow, not feature tables.
		\end{itemize}
	\end{frame}
	
	% ---------------------------------------------------------------
	\begin{frame}{Question 4}
		Which method creates a feature table with \texttt{FeatureEngineeringClient}?
		\begin{enumerate}[label=\Alph*.]
			\item \texttt{fe.create\_table()}
			\item \texttt{fe.write\_table()}
			\item \texttt{fe.create\_training\_set()}
			\item \texttt{fe.log\_model()}
		\end{enumerate}
	\end{frame}
	\begin{frame}{Answer 4}
		\textbf{Correct: A} --- \texttt{create\_table()} creates the table.\\
		\vspace{0.5em}
		\textbf{Wrong:}
		\begin{itemize}
			\item B writes features to an existing table.
			\item C builds a training set.
			\item D logs a model.
		\end{itemize}
	\end{frame}
	
	% ---------------------------------------------------------------
	\begin{frame}{Question 5}
		Which method writes features to an existing feature table?
		\begin{enumerate}[label=\Alph*.]
			\item \texttt{fe.create\_table()}
			\item \texttt{fe.write\_table()}
			\item \texttt{fe.score\_batch()}
			\item \texttt{fe.get\_table()}
		\end{enumerate}
	\end{frame}
	\begin{frame}{Answer 5}
		\textbf{Correct: B} --- \texttt{write\_table()} writes features.\\
		\vspace{0.5em}
		\textbf{Wrong:}
		\begin{itemize}
			\item A creates a table.
			\item C scores a batch.
			\item D retrieves a table.
		\end{itemize}
	\end{frame}
	
	% ---------------------------------------------------------------
	\begin{frame}{Question 6}
		What does \texttt{fe.create\_training\_set()} provide that a plain DataFrame join does not?
		\begin{enumerate}[label=\Alph*.]
			\item Automatic one-hot encoding
			\item Point-in-time lookups aligned to the label timestamp
			\item Model registration
			\item Hyperparameter tuning
		\end{enumerate}
	\end{frame}
	\begin{frame}{Answer 6}
		\textbf{Correct: B} --- training sets use point-in-time lookups aligned to label timestamps.\\
		\vspace{0.5em}
		\textbf{Wrong:}
		\begin{itemize}
			\item A is a separate preprocessing step.
			\item C is done by \texttt{register\_model}.
			\item D is done by Hyperopt.
		\end{itemize}
	\end{frame}
	
	% ---------------------------------------------------------------
	\begin{frame}{Question 7}
		A model performs well offline but degrades in production. Review shows a feature was computed using data available only after the label outcome. What is the diagnosis?
		\begin{enumerate}[label=\Alph*.]
			\item Underfitting
			\item Training-serving skew / target leakage
			\item Class imbalance
			\item Concept drift
		\end{enumerate}
	\end{frame}
	\begin{frame}{Answer 7}
		\textbf{Correct: B} --- post-label data causes leakage and skew.\\
		\vspace{0.5em}
		\textbf{Wrong:}
		\begin{itemize}
			\item A would show poor training performance.
			\item C does not explain post-label features.
			\item D is a change over time, not a data leak.
		\end{itemize}
	\end{frame}
	
	% ---------------------------------------------------------------
	\begin{frame}{Question 8}
		What is the difference between online and offline feature tables?
		\begin{enumerate}[label=\Alph*.]
			\item Online is for training; offline is for serving
			\item Online serves low-latency real-time inference; offline is for batch training and scoring
			\item They are identical
			\item Offline is deprecated
		\end{enumerate}
	\end{frame}
	\begin{frame}{Answer 8}
		\textbf{Correct: B} --- online tables serve low-latency inference; offline tables are for batch.\\
		\vspace{0.5em}
		\textbf{Wrong:}
		\begin{itemize}
			\item A reverses the roles.
			\item C is false.
			\item D is false; offline tables are not deprecated.
		\end{itemize}
	\end{frame}
	
	% ---------------------------------------------------------------
	\begin{frame}{Question 9}
		A feature-aware model was logged with \texttt{fs.log\_model()} and a \texttt{batch\_df} contains only the primary key \texttt{customer\_id}. Which call generates predictions by automatically joining the required features?
		\begin{enumerate}[label=\Alph*.]
			\item \texttt{model.predict(batch\_df)}
			\item \texttt{fs.score\_batch(model\_uri, batch\_df)}
			\item \texttt{mlflow.pyfunc.load\_model(model\_uri).predict(batch\_df)}
			\item \texttt{batch\_df.join(feature\_table)}
		\end{enumerate}
	\end{frame}
	\begin{frame}{Answer 9}
		\textbf{Correct: B} --- \texttt{fs.score\_batch()} uses logged feature metadata to join features automatically.\\
		\vspace{0.5em}
		\textbf{Wrong:}
		\begin{itemize}
			\item A assumes the features are already present.
			\item C also assumes features are in the DataFrame.
			\item D is a plain join without automatic lookup.
		\end{itemize}
	\end{frame}
	
	% ---------------------------------------------------------------
	\begin{frame}{Question 10}
		Which namespace format is correct for a feature table in Unity Catalog?
		\begin{enumerate}[label=\Alph*.]
			\item \texttt{catalog.schema.table}
			\item \texttt{database.table}
			\item \texttt{workspace.table}
			\item \texttt{hive\_metastore.table}
		\end{enumerate}
	\end{frame}
	\begin{frame}{Answer 10}
		\textbf{Correct: A} --- Unity Catalog uses \texttt{catalog.schema.table}.\\
		\vspace{0.5em}
		\textbf{Wrong:}
		\begin{itemize}
			\item B is a two-level format.
			\item C is not a Unity Catalog format.
			\item D is a legacy Hive reference.
		\end{itemize}
	\end{frame}
	
	% ---------------------------------------------------------------
	\begin{frame}{Question 11}
		Which MLflow API call manually logs a metric?
		\begin{enumerate}[label=\Alph*.]
			\item \texttt{mlflow.log\_metric()}
			\item \texttt{mlflow.log\_param()}
			\item \texttt{mlflow.log\_artifact()}
			\item \texttt{mlflow.set\_tag()}
		\end{enumerate}
	\end{frame}
	\begin{frame}{Answer 11}
		\textbf{Correct: A} --- \texttt{log\_metric()} logs metrics.\\
		\vspace{0.5em}
		\textbf{Wrong:}
		\begin{itemize}
			\item B logs parameters.
			\item C logs artifacts.
			\item D sets tags.
		\end{itemize}
	\end{frame}
	
	% ---------------------------------------------------------------
	\begin{frame}{Question 12}
		Which MLflow API call registers a model in Unity Catalog?
		\begin{enumerate}[label=\Alph*.]
			\item \texttt{mlflow.log\_model()}
			\item \texttt{mlflow.register\_model()}
			\item \texttt{mlflow.save\_model()}
			\item \texttt{mlflow.start\_run()}
		\end{enumerate}
	\end{frame}
	\begin{frame}{Answer 12}
		\textbf{Correct: B} --- \texttt{register\_model()} registers a model.\\
		\vspace{0.5em}
		\textbf{Wrong:}
		\begin{itemize}
			\item A logs a model artifact.
			\item C saves locally.
			\item D starts a run.
		\end{itemize}
	\end{frame}
	
	% ---------------------------------------------------------------
	\begin{frame}{Question 13}
		What replaced the old Staging and Production stages for model promotion in Unity Catalog?
		\begin{enumerate}[label=\Alph*.]
			\item Tags only
			\item Aliases such as \texttt{@champion}
			\item Notebook revisions
			\item Git branches
		\end{enumerate}
	\end{frame}
	\begin{frame}{Answer 13}
		\textbf{Correct: B} --- aliases replace stages.\\
		\vspace{0.5em}
		\textbf{Wrong:}
		\begin{itemize}
			\item A tags do not replace promotion.
			\item C tracks code, not model stages.
			\item D tracks code versions.
		\end{itemize}
	\end{frame}
	
	% ---------------------------------------------------------------
	\begin{frame}{Question 14}
		How do you promote a challenger model to champion using aliases?
		\begin{enumerate}[label=\Alph*.]
			\item Delete the old model version
			\item Move the \texttt{champion} alias to the new model version
			\item Rename the model
			\item Create a new experiment
		\end{enumerate}
	\end{frame}
	\begin{frame}{Answer 14}
		\textbf{Correct: B} --- promotion is moving the alias.\\
		\vspace{0.5em}
		\textbf{Wrong:}
		\begin{itemize}
			\item A deletes history.
			\item C does not change deployment.
			\item D is unrelated.
		\end{itemize}
	\end{frame}
	
	% ---------------------------------------------------------------
	\begin{frame}{Question 15}
		What does \texttt{mlflow.search\_runs()} return?
		\begin{enumerate}[label=\Alph*.]
			\item A single run ID
			\item A pandas DataFrame of runs that can be sorted and filtered
			\item A trained model
			\item A list of experiments
		\end{enumerate}
	\end{frame}
	\begin{frame}{Answer 15}
		\textbf{Correct: B} --- it returns a pandas DataFrame of runs.\\
		\vspace{0.5em}
		\textbf{Wrong:}
		\begin{itemize}
			\item A is too narrow.
			\item C is a model object.
			\item D lists experiments, not runs.
		\end{itemize}
	\end{frame}
	
	% ---------------------------------------------------------------
	\begin{frame}{Question 16}
		A run ID is stored in \texttt{run\_id} and the logged model name is \texttt{"model"}. Which line registers it as \texttt{"best\_model"}?
		\begin{enumerate}[label=\Alph*.]
			\item \texttt{mlflow.register\_model(run\_id, "best\_model")}
			\item \texttt{mlflow.register\_model(f"runs:/\{run\_id\}/model", "best\_model")}
			\item \texttt{mlflow.register\_model(f"runs:/\{run\_id\}/model")}
			\item \texttt{mlflow.register\_model(f"runs:/\{run\_id\}/best\_model", "model")}
		\end{enumerate}
	\end{frame}
	\begin{frame}{Answer 16}
		\textbf{Correct: B} --- the URI must include \texttt{runs:/} and the artifact path.\\
		\vspace{0.5em}
		\textbf{Wrong:}
		\begin{itemize}
			\item A passes only the run ID, not a model URI.
			\item C omits the registry name.
			\item D swaps artifact path and name.
		\end{itemize}
	\end{frame}
	
	% ---------------------------------------------------------------
	\begin{frame}{Question 17}
		A data scientist wants one parent run for tuning and a child run per trial, starting runs manually. What should they do?
		\begin{enumerate}[label=\Alph*.]
			\item Turn on Databricks Autologging
			\item Specify \texttt{nested=True} when starting each child run
			\item Start each child run inside the parent's code block without \texttt{nested}
			\item Give child runs the same experiment ID
		\end{enumerate}
	\end{frame}
	\begin{frame}{Answer 17}
		\textbf{Correct: B} --- child runs must be started with \texttt{nested=True}.\\
		\vspace{0.5em}
		\textbf{Wrong:}
		\begin{itemize}
			\item A does not create nesting.
			\item C raises an error.
			\item D only groups runs in the same experiment.
		\end{itemize}
	\end{frame}
	
	% ---------------------------------------------------------------
	\begin{frame}{Question 18}
		Which UI element on the MLflow experiment page opens the notebook that created a run?
		\begin{enumerate}[label=\Alph*.]
			\item The \texttt{MLmodel} artifact
			\item The \texttt{Models} link
			\item The \texttt{Source} link
			\item The \texttt{Start Time} link
		\end{enumerate}
	\end{frame}
	\begin{frame}{Answer 18}
		\textbf{Correct: C} --- the Source link opens the notebook snapshot.\\
		\vspace{0.5em}
		\textbf{Wrong:}
		\begin{itemize}
			\item A is model metadata.
			\item B links to the registry.
			\item D opens run details.
		\end{itemize}
	\end{frame}
	
	% ---------------------------------------------------------------
	\begin{frame}{Question 19}
		Which function searches runs ordered by a metric ascending to find the lowest RMSE?
		\begin{enumerate}[label=\Alph*.]
			\item \texttt{mlflow.search\_runs(experiment\_id, order\_by=["metrics.rmse DESC"])}
			\item \texttt{mlflow.best\_run(experiment\_id, order\_by=["metrics.rmse"])}
			\item \texttt{mlflow.search\_runs(experiment\_id, order\_by=["metrics.rmse ASC"])}
			\item \texttt{mlflow.search\_runs(experiment\_id, filter="rmse < 1")}
		\end{enumerate}
	\end{frame}
	\begin{frame}{Answer 19}
		\textbf{Correct: C} --- ascending order puts the lowest RMSE first.\\
		\vspace{0.5em}
		\textbf{Wrong:}
		\begin{itemize}
			\item A orders descending.
			\item B is not the API.
			\item D filters rather than sorts.
		\end{itemize}
	\end{frame}
	
	% ---------------------------------------------------------------
	\begin{frame}{Question 20}
		On which page do you transition a registered model version from Staging to Production?
		\begin{enumerate}[label=\Alph*.]
			\item The MLflow Model Registry home page
			\item The experiment page
			\item The model version page
			\item The model page only
		\end{enumerate}
	\end{frame}
	\begin{frame}{Answer 20}
		\textbf{Correct: C} --- stage transitions happen on the model version page.\\
		\vspace{0.5em}
		\textbf{Wrong:}
		\begin{itemize}
			\item A lists all models.
			\item B shows runs.
			\item D lists versions but not the direct controls.
		\end{itemize}
	\end{frame}
	
	% ---------------------------------------------------------------
	\begin{frame}{Question 21}
		Which MLflow API call logs an artifact?
		\begin{enumerate}[label=\Alph*.]
			\item \texttt{mlflow.log\_metric()}
			\item \texttt{mlflow.log\_param()}
			\item \texttt{mlflow.log\_artifact()}
			\item \texttt{mlflow.set\_tag()}
		\end{enumerate}
	\end{frame}
	\begin{frame}{Answer 21}
		\textbf{Correct: C} --- \texttt{log\_artifact()} logs files.\\
		\vspace{0.5em}
		\textbf{Wrong:}
		\begin{itemize}
			\item A logs metrics.
			\item B logs parameters.
			\item D sets tags.
		\end{itemize}
	\end{frame}
	
	% ---------------------------------------------------------------
	\begin{frame}{Question 22}
		What is the simplest Databricks-native way to make every training run reproducible months later?
		\begin{enumerate}[label=\Alph*.]
			\item Save notebooks to a local drive
			\item Use MLflow tracking to log parameters, metrics, artifacts, and the model
			\item Print metrics to stdout
			\item Rely on cluster logs
		\end{enumerate}
	\end{frame}
	\begin{frame}{Answer 22}
		\textbf{Correct: B} --- MLflow tracking logs all run data for reproducibility.\\
		\vspace{0.5em}
		\textbf{Wrong:}
		\begin{itemize}
			\item A is not centralized.
			\item C is not durable.
			\item D may be ephemeral.
		\end{itemize}
	\end{frame}
	
	% ---------------------------------------------------------------
	\begin{frame}{Question 23}
		What does Databricks AutoML produce besides a trained model?
		\begin{enumerate}[label=\Alph*.]
			\item Only a model artifact
			\item A leaderboard, editable notebooks, and an automated EDA notebook
			\item A production endpoint
			\item A data warehouse
		\end{enumerate}
	\end{frame}
	\begin{frame}{Answer 23}
		\textbf{Correct: B} --- AutoML returns a leaderboard, notebooks, and EDA.\\
		\vspace{0.5em}
		\textbf{Wrong:}
		\begin{itemize}
			\item A is incomplete.
			\item C is a deployment step.
			\item D is not produced.
		\end{itemize}
	\end{frame}
	
	% ---------------------------------------------------------------
	\begin{frame}{Question 24}
		Which API runs classification AutoML?
		\begin{enumerate}[label=\Alph*.]
			\item \texttt{databricks.automl.classify()}
			\item \texttt{databricks.automl.regress()}
			\item \texttt{databricks.automl.forecast()}
			\item \texttt{databricks.automl.train()}
		\end{enumerate}
	\end{frame}
	\begin{frame}{Answer 24}
		\textbf{Correct: A} --- \texttt{classify()} runs classification AutoML.\\
		\vspace{0.5em}
		\textbf{Wrong:}
		\begin{itemize}
			\item B is regression.
			\item C is forecasting.
			\item D is not an API.
		\end{itemize}
	\end{frame}
	
	% ---------------------------------------------------------------
	\begin{frame}{Question 25}
		AutoML is least appropriate when:
		\begin{enumerate}[label=\Alph*.]
			\item You need a quick baseline
			\item You need a highly customized deep learning architecture
			\item You want automatic feature engineering
			\item You want to compare multiple algorithms
		\end{enumerate}
	\end{frame}
	\begin{frame}{Answer 25}
		\textbf{Correct: B} --- AutoML is less suited to highly custom deep learning.\\
		\vspace{0.5em}
		\textbf{Wrong:}
		\begin{itemize}
			\item A is a strength of AutoML.
			\item C is a strength of AutoML.
			\item D is a strength of AutoML.
		\end{itemize}
	\end{frame}
	
	% ---------------------------------------------------------------
	\begin{frame}{Question 26}
		Which step must a data scientist perform outside an AutoML experiment?
		\begin{enumerate}[label=\Alph*.]
			\item Model tuning
			\item Model evaluation
			\item Model deployment
			\item Exploratory data analysis
		\end{enumerate}
	\end{frame}
	\begin{frame}{Answer 26}
		\textbf{Correct: C} --- deployment is outside AutoML.\\
		\vspace{0.5em}
		\textbf{Wrong:}
		\begin{itemize}
			\item A is automated.
			\item B is automated.
			\item D is automated via the EDA notebook.
		\end{itemize}
	\end{frame}
	
	% ---------------------------------------------------------------
	\begin{frame}{Question 27}
		How do you begin using Databricks Runtime ML?
		\begin{enumerate}[label=\Alph*.]
			\item Add an init script line
			\item Check a Databricks Runtime ML box
			\item Select a Runtime ML version from the Databricks Runtime Version dropdown
			\item Set \texttt{runtime-version="ml"} in the Spark session
		\end{enumerate}
	\end{frame}
	\begin{frame}{Answer 27}
		\textbf{Correct: C} --- Runtime ML is selected from the runtime dropdown.\\
		\vspace{0.5em}
		\textbf{Wrong:}
		\begin{itemize}
			\item A does not select the runtime.
			\item B there is no such box.
			\item D does not change the runtime.
		\end{itemize}
	\end{frame}
	
	% ---------------------------------------------------------------
	\begin{frame}{Question 28}
		Why is Databricks Runtime ML convenient for MLflow users?
		\begin{enumerate}[label=\Alph*.]
			\item It removes the need for a cluster
			\item It preinstalls MLflow and other ML libraries
			\item It disables logging
			\item It only supports R
		\end{enumerate}
	\end{frame}
	\begin{frame}{Answer 28}
		\textbf{Correct: B} --- Runtime ML preinstalls MLflow and ML libraries.\\
		\vspace{0.5em}
		\textbf{Wrong:}
		\begin{itemize}
			\item A is false.
			\item C is false.
			\item D is false; it supports Python and more.
		\end{itemize}
	\end{frame}
	
	% ---------------------------------------------------------------
	\begin{frame}{Question 29}
		Which evaluation metric is not suitable for AutoML regression experiments?
		\begin{enumerate}[label=\Alph*.]
			\item F1
			\item R-squared
			\item MAE
			\item MSE
		\end{enumerate}
	\end{frame}
	\begin{frame}{Answer 29}
		\textbf{Correct: A} --- F1 is for classification, not regression.\\
		\vspace{0.5em}
		\textbf{Wrong:}
		\begin{itemize}
			\item B is a regression metric.
			\item C is a regression metric.
			\item D is a regression metric.
		\end{itemize}
	\end{frame}
	
	% ---------------------------------------------------------------
	\begin{frame}{Question 30}
		Which file typically defines a model's format and dependencies in an MLflow run?
		\begin{enumerate}[label=\Alph*.]
			\item \texttt{requirements.txt}
			\item \texttt{MLmodel}
			\item \texttt{README.md}
			\item \texttt{conda.yaml} only
		\end{enumerate}
	\end{frame}
	\begin{frame}{Answer 30}
		\textbf{Correct: B} --- \texttt{MLmodel} defines format and dependencies.\\
		\vspace{0.5em}
		\textbf{Wrong:}
		\begin{itemize}
			\item A is a dependency list.
			\item C is documentation.
			\item D is one dependency spec, not the format file.
		\end{itemize}
	\end{frame}
	
	% ---------------------------------------------------------------
	\begin{frame}{Question 31}
		What does \texttt{fs.log\_model()} add to a model artifact?
		\begin{enumerate}[label=\Alph*.]
			\item Hyperparameters
			\item Feature lookup metadata for automatic joins at inference
			\item A confusion matrix
			\item A serving endpoint
		\end{enumerate}
	\end{frame}
	\begin{frame}{Answer 31}
		\textbf{Correct: B} --- it packages feature lookup metadata.\\
		\vspace{0.5em}
		\textbf{Wrong:}
		\begin{itemize}
			\item A is logged separately.
			\item C is not added by this call.
			\item D is created later.
		\end{itemize}
	\end{frame}
	
	% ---------------------------------------------------------------
	\begin{frame}{Question 32}
		Which is a benefit of registering models in Unity Catalog over the workspace registry?
		\begin{enumerate}[label=\Alph*.]
			\item Faster training
			\item Lineage, governance, and cross-workspace access
			\item Automatic hyperparameter tuning
			\item Free storage
		\end{enumerate}
	\end{frame}
	\begin{frame}{Answer 32}
		\textbf{Correct: B} --- Unity Catalog adds lineage, governance, and access.\\
		\vspace{0.5em}
		\textbf{Wrong:}
		\begin{itemize}
			\item A is false.
			\item C is false.
			\item D is false.
		\end{itemize}
	\end{frame}
	
	% ---------------------------------------------------------------
	\begin{frame}{Question 33}
		What does an MLflow experiment group together?
		\begin{enumerate}[label=\Alph*.]
			\item Runs of a related ML task
			\item Clusters
			\item Notebooks only
			\item Feature tables
		\end{enumerate}
	\end{frame}
	\begin{frame}{Answer 33}
		\textbf{Correct: A} --- experiments group related runs.\\
		\vspace{0.5em}
		\textbf{Wrong:}
		\begin{itemize}
			\item B clusters are compute.
			\item C notebooks are source code.
			\item D feature tables are data assets.
		\end{itemize}
	\end{frame}
	
	% ---------------------------------------------------------------
	\begin{frame}{Question 34}
		Which API call starts an MLflow run?
		\begin{enumerate}[label=\Alph*.]
			\item \texttt{mlflow.start\_run()}
			\item \texttt{mlflow.log\_run()}
			\item \texttt{mlflow.open\_run()}
			\item \texttt{mlflow.new\_run()}
		\end{enumerate}
	\end{frame}
	\begin{frame}{Answer 34}
		\textbf{Correct: A} --- \texttt{start\_run()} starts a run.\\
		\vspace{0.5em}
		\textbf{Wrong:}
		\begin{itemize}
			\item B is not an API.
			\item C is not an API.
			\item D is not an API.
		\end{itemize}
	\end{frame}
	
	% ---------------------------------------------------------------
	\begin{frame}{Question 35}
		Which is true about aliases in the Unity Catalog model registry?
		\begin{enumerate}[label=\Alph*.]
			\item An alias can point to many versions at once
			\item An alias points to one model version and can be moved
			\item Aliases are immutable
			\item Aliases replace experiments
		\end{enumerate}
	\end{frame}
	\begin{frame}{Answer 35}
		\textbf{Correct: B} --- an alias points to one version and can be moved.\\
		\vspace{0.5em}
		\textbf{Wrong:}
		\begin{itemize}
			\item A is false.
			\item C aliases are movable.
			\item D aliases are model-level, not experiment-level.
		\end{itemize}
	\end{frame}
	
	% ---------------------------------------------------------------
	\begin{frame}{Question 36}
		Which object stores the actual model files and metadata for a run?
		\begin{enumerate}[label=\Alph*.]
			\item The artifact store
			\item The parameter log
			\item The metric log
			\item The experiment tag
		\end{enumerate}
	\end{frame}
	\begin{frame}{Answer 36}
		\textbf{Correct: A} --- artifacts live in the artifact store.\\
		\vspace{0.5em}
		\textbf{Wrong:}
		\begin{itemize}
			\item B stores parameters.
			\item C stores metrics.
			\item D stores tags.
		\end{itemize}
	\end{frame}
	
	% ---------------------------------------------------------------
	\begin{frame}{Question 37}
		Why should secrets not be logged as MLflow parameters?
		\begin{enumerate}[label=\Alph*.]
			\item Parameters are stored in plain text and visible to experiment users
			\item Parameters are not persisted
			\item Parameters are encrypted by default
			\item Parameters are deleted after the run
		\end{enumerate}
	\end{frame}
	\begin{frame}{Answer 37}
		\textbf{Correct: A} --- parameters are stored in plain text.\\
		\vspace{0.5em}
		\textbf{Wrong:}
		\begin{itemize}
			\item B is false; they are persisted.
			\item C is false.
			\item D is false.
		\end{itemize}
	\end{frame}
	
	% ---------------------------------------------------------------
	\begin{frame}{Question 38}
		Which is a primary benefit of using MLflow Model Registry?
		\begin{enumerate}[label=\Alph*.]
			\item It replaces all experiments
			\item It manages model versions, aliases, and lifecycle
			\item It trains models
			\item It stores feature tables
		\end{enumerate}
	\end{frame}
	\begin{frame}{Answer 38}
		\textbf{Correct: B} --- the registry manages versions, aliases, and lifecycle.\\
		\vspace{0.5em}
		\textbf{Wrong:}
		\begin{itemize}
			\item A experiments still exist.
			\item C it does not train models.
			\item D feature tables are in Feature Store.
		\end{itemize}
	\end{frame}
	
	% ---------------------------------------------------------------
	\begin{frame}{Question 39}
		A Spark DataFrame \texttt{spark\_df} must keep only rows where \texttt{price > 0}. Which code block works?
		\begin{enumerate}[label=\Alph*.]
			\item \texttt{spark\_df[spark\_df["price"] > 0]}
			\item \texttt{spark\_df.filter(col("price") > 0)}
			\item \texttt{SELECT * FROM spark\_df WHERE price > 0}
			\item \texttt{spark\_df.loc[spark\_df["price"] > 0,:]}
		\end{enumerate}
	\end{frame}
	\begin{frame}{Answer 39}
		\textbf{Correct: B} --- \texttt{filter(col(...))} is the PySpark way.\\
		\vspace{0.5em}
		\textbf{Wrong:}
		\begin{itemize}
			\item A is pandas-like and not PySpark-idiomatic.
			\item C is SQL, not DataFrame code.
			\item D is pandas syntax.
		\end{itemize}
	\end{frame}
	
	% ---------------------------------------------------------------
	\begin{frame}{Question 40}
		Which line returns count, mean, standard deviation, min, max, and quantiles for numeric columns?
		\begin{enumerate}[label=\Alph*.]
			\item \texttt{spark\_df.summary()}
			\item \texttt{spark\_df.stats()}
			\item \texttt{spark\_df.describe().head()}
			\item \texttt{spark\_df.printSchema()}
		\end{enumerate}
	\end{frame}
	\begin{frame}{Answer 40}
		\textbf{Correct: A} --- \texttt{summary()} returns count, mean, std, min, max, and quantiles.\\
		\vspace{0.5em}
		\textbf{Wrong:}
		\begin{itemize}
			\item B is not an API.
			\item C lacks quantiles.
			\item D prints the schema.
		\end{itemize}
	\end{frame}
	
	% ---------------------------------------------------------------
	\begin{frame}{Question 41}
		How do you remove outliers using IQR in Spark?
		\begin{enumerate}[label=\Alph*.]
			\item Drop the top 10\% of rows
			\item Filter values outside $[Q1 - 1.5 \cdot IQR,\; Q3 + 1.5 \cdot IQR]$
			\item Filter values outside $[Q1 - IQR,\; Q3 + IQR]$
			\item Replace all values with the mean
		\end{enumerate}
	\end{frame}
	\begin{frame}{Answer 41}
		\textbf{Correct: B} --- the $1.5 \times IQR$ rule is standard.\\
		\vspace{0.5em}
		\textbf{Wrong:}
		\begin{itemize}
			\item A is not IQR-based.
			\item C uses a wrong threshold.
			\item D does not remove outliers.
		\end{itemize}
	\end{frame}
	
	% ---------------------------------------------------------------
	\begin{frame}{Question 42}
		When is median imputation preferred over mean imputation?
		\begin{enumerate}[label=\Alph*.]
			\item When data is symmetric
			\item When data is skewed or contains extreme outliers
			\item When there are no missing values
			\item Always
		\end{enumerate}
	\end{frame}
	\begin{frame}{Answer 42}
		\textbf{Correct: B} --- median is robust to skew and outliers.\\
		\vspace{0.5em}
		\textbf{Wrong:}
		\begin{itemize}
			\item A mean is fine for symmetric data.
			\item C no imputation is needed.
			\item D is false.
		\end{itemize}
	\end{frame}
	
	% ---------------------------------------------------------------
	\begin{frame}{Question 43}
		A data scientist replaces missing values with the median. A colleague says valuable information is lost. Which approach preserves the most information?
		\begin{enumerate}[label=\Alph*.]
			\item Impute with the mean instead
			\item Let the algorithm handle missing values
			\item Drop all features with missing values
			\item Add a binary indicator feature marking which values were imputed
		\end{enumerate}
	\end{frame}
	\begin{frame}{Answer 43}
		\textbf{Correct: D} --- indicator features preserve missingness information.\\
		\vspace{0.5em}
		\textbf{Wrong:}
		\begin{itemize}
			\item A still loses missingness information.
			\item B not all algorithms handle missing values well.
			\item C discards features entirely.
		\end{itemize}
	\end{frame}
	
	% ---------------------------------------------------------------
	\begin{frame}{Question 44}
		One-hot encoding is least appropriate for:
		\begin{enumerate}[label=\Alph*.]
			\item Nominal categorical features
			\item Very high-cardinality categorical features
			\item Low-cardinality categorical features
			\item Tree-based models
		\end{enumerate}
	\end{frame}
	\begin{frame}{Answer 44}
		\textbf{Correct: B} --- high cardinality causes dimensionality explosion.\\
		\vspace{0.5em}
		\textbf{Wrong:}
		\begin{itemize}
			\item A is a typical use case.
			\item C is a typical use case.
			\item D trees can use one-hot but are not the issue.
		\end{itemize}
	\end{frame}
	
	% ---------------------------------------------------------------
	\begin{frame}{Question 45}
		A feature repository one-hot encodes all categorical variables. Why is this problematic?
		\begin{enumerate}[label=\Alph*.]
			\item One-hot encoding is not supported by most libraries
			\item It depends on target values
			\item It is computationally intensive and should be done on small samples
			\item It is a potentially problematic strategy for some ML algorithms due to dimensionality and sparsity
		\end{enumerate}
	\end{frame}
	\begin{frame}{Answer 45}
		\textbf{Correct: D} --- premature one-hot can hurt some algorithms.\\
		\vspace{0.5em}
		\textbf{Wrong:}
		\begin{itemize}
			\item A is false.
			\item B one-hot does not use target values.
			\item C is not the main issue.
		\end{itemize}
	\end{frame}
	
	% ---------------------------------------------------------------
	\begin{frame}{Question 46}
		How do you one-hot encode in Spark ML?
		\begin{enumerate}[label=\Alph*.]
			\item \texttt{OneHotEncoder} alone
			\item \texttt{StringIndexer} then \texttt{OneHotEncoder}
			\item \texttt{VectorAssembler} then \texttt{StringIndexer}
			\item \texttt{Imputer} then \texttt{OneHotEncoder}
		\end{enumerate}
	\end{frame}
	\begin{frame}{Answer 46}
		\textbf{Correct: B} --- Spark ML uses StringIndexer then OneHotEncoder.\\
		\vspace{0.5em}
		\textbf{Wrong:}
		\begin{itemize}
			\item A needs indexed input first.
			\item C is wrong order.
			\item D imputation is unrelated.
		\end{itemize}
	\end{frame}
	
	% ---------------------------------------------------------------
	\begin{frame}{Question 47}
		What does \texttt{VectorAssembler} do?
		\begin{enumerate}[label=\Alph*.]
			\item Combines multiple columns into a single feature vector
			\item Splits a vector into columns
			\item Encodes strings
			\item Imputes missing values
		\end{enumerate}
	\end{frame}
	\begin{frame}{Answer 47}
		\textbf{Correct: A} --- VectorAssembler combines columns into a vector.\\
		\vspace{0.5em}
		\textbf{Wrong:}
		\begin{itemize}
			\item B is the opposite operation.
			\item C is StringIndexer.
			\item D is Imputer.
		\end{itemize}
	\end{frame}
	
	% ---------------------------------------------------------------
	\begin{frame}{Question 48}
		A data scientist creates an \texttt{Imputer} and calls \texttt{transform} directly on it. Why does it fail?
		\begin{enumerate}[label=\Alph*.]
			\item It does not impute both training and test sets
			\item \texttt{inputCols} and \texttt{outputCols} must be identical
			\item It should call \texttt{fit} instead of \texttt{transform}
			\item It does not fit the imputer to create an \texttt{ImputerModel}
		\end{enumerate}
	\end{frame}
	\begin{frame}{Answer 48}
		\textbf{Correct: D} --- the imputer must be fit first.\\
		\vspace{0.5em}
		\textbf{Wrong:}
		\begin{itemize}
			\item A is not the immediate cause.
			\item B is false; columns can differ.
			\item C is imprecise; you fit then transform.
		\end{itemize}
	\end{frame}
	
	% ---------------------------------------------------------------
	\begin{frame}{Question 49}
		When is log transformation appropriate?
		\begin{enumerate}[label=\Alph*.]
			\item For right-skewed positive features
			\item For left-skewed negative features
			\item For categorical features
			\item For binary features
		\end{enumerate}
	\end{frame}
	\begin{frame}{Answer 49}
		\textbf{Correct: A} --- log transform suits right-skewed positive data.\\
		\vspace{0.5em}
		\textbf{Wrong:}
		\begin{itemize}
			\item B is not a typical use case.
			\item C is for categorical data.
			\item D is for binary data.
		\end{itemize}
	\end{frame}
	
	% ---------------------------------------------------------------
	\begin{frame}{Question 50}
		A pandas notebook must scale to big data with the least refactoring effort. What approach is recommended?
		\begin{enumerate}[label=\Alph*.]
			\item Refactor to process data in parallel
			\item Refactor to the PySpark DataFrame API
			\item Refactor to the Scala Dataset API
			\item Refactor to Spark SQL
			\item Refactor to the pandas API on Spark
		\end{enumerate}
	\end{frame}
	\begin{frame}{Answer 50}
		\textbf{Correct: E} --- pandas API on Spark minimizes code changes.\\
		\vspace{0.5em}
		\textbf{Wrong:}
		\begin{itemize}
			\item A--D require larger rewrites.
		\end{itemize}
	\end{frame}
	
	% ---------------------------------------------------------------
	\begin{frame}{Question 51}
		Which function applies a Python function to each group of a grouped DataFrame and returns a DataFrame?
		\begin{enumerate}[label=\Alph*.]
			\item \texttt{applyInPandas}
			\item \texttt{groupedApplyInPandas}
			\item \texttt{mapInPandas}
			\item \texttt{predict}
		\end{enumerate}
	\end{frame}
	\begin{frame}{Answer 51}
		\textbf{Correct: A} --- \texttt{applyInPandas} applies per group and returns a DataFrame.\\
		\vspace{0.5em}
		\textbf{Wrong:}
		\begin{itemize}
			\item B is not an API.
			\item C maps partitions, not groups.
			\item D is not an API.
		\end{itemize}
	\end{frame}
	
	% ---------------------------------------------------------------
	\begin{frame}{Question 52}
		A data scientist wants group-specific model inference in parallel. Which call fills the blank in \texttt{df.groupby("device\_id").\_\_\_(apply\_model, schema=...) }?
		\begin{enumerate}[label=\Alph*.]
			\item \texttt{applyInPandas}
			\item \texttt{groupedApplyInPandas}
			\item \texttt{mapInPandas}
			\item \texttt{predict}
		\end{enumerate}
	\end{frame}
	\begin{frame}{Answer 52}
		\textbf{Correct: A} --- \texttt{applyInPandas} is the grouped API.\\
		\vspace{0.5em}
		\textbf{Wrong:}
		\begin{itemize}
			\item B is not an API.
			\item C maps partitions.
			\item D is not an API.
		\end{itemize}
	\end{frame}
	
	% ---------------------------------------------------------------
	\begin{frame}{Question 53}
		Which imputation strategy is most appropriate for a numeric feature with extreme outliers?
		\begin{enumerate}[label=\Alph*.]
			\item Mean
			\item Median
			\item Mode
			\item Zero
		\end{enumerate}
	\end{frame}
	\begin{frame}{Answer 53}
		\textbf{Correct: B} --- median is robust to outliers.\\
		\vspace{0.5em}
		\textbf{Wrong:}
		\begin{itemize}
			\item A is distorted by outliers.
			\item C is for categorical data.
			\item D can bias the distribution.
		\end{itemize}
	\end{frame}
	
	% ---------------------------------------------------------------
	\begin{frame}{Question 54}
		What does a boxplot show?
		\begin{enumerate}[label=\Alph*.]
			\item Only the mean
			\item Median, quartiles, and outliers
			\item Correlation
			\item Class balance
		\end{enumerate}
	\end{frame}
	\begin{frame}{Answer 54}
		\textbf{Correct: B} --- boxplots show median, quartiles, and outliers.\\
		\vspace{0.5em}
		\textbf{Wrong:}
		\begin{itemize}
			\item A is incomplete.
			\item C is a scatter/heatmap feature.
			\item D is a bar chart feature.
		\end{itemize}
	\end{frame}
	
	% ---------------------------------------------------------------
	\begin{frame}{Question 55}
		Which technique helps detect multicollinearity?
		\begin{enumerate}[label=\Alph*.]
			\item Variance Inflation Factor
			\item Silhouette score
			\item ROC AUC
			\item Confusion matrix
		\end{enumerate}
	\end{frame}
	\begin{frame}{Answer 55}
		\textbf{Correct: A} --- VIF detects multicollinearity.\\
		\vspace{0.5em}
		\textbf{Wrong:}
		\begin{itemize}
			\item B measures clustering quality.
			\item C measures classification ranking.
			\item D evaluates classification outcomes.
		\end{itemize}
	\end{frame}
	
	% ---------------------------------------------------------------
	\begin{frame}{Question 56}
		Why is stratification important for classification splits?
		\begin{enumerate}[label=\Alph*.]
			\item It preserves class proportions in each split
			\item It increases training speed
			\item It removes outliers
			\item It ensures reproducibility only
		\end{enumerate}
	\end{frame}
	\begin{frame}{Answer 56}
		\textbf{Correct: A} --- stratification preserves class proportions.\\
		\vspace{0.5em}
		\textbf{Wrong:}
		\begin{itemize}
			\item B is not its purpose.
			\item C is a different preprocessing step.
			\item D is not its main purpose.
		\end{itemize}
	\end{frame}
	
	% ---------------------------------------------------------------
	\begin{frame}{Question 57}
		Which is a symptom of an imbalanced dataset?
		\begin{enumerate}[label=\Alph*.]
			\item A 50/50 class split
			\item A 99\% negative / 1\% positive split
			\item A continuous target with normal distribution
			\item Missing values
		\end{enumerate}
	\end{frame}
	\begin{frame}{Answer 57}
		\textbf{Correct: B} --- 99/1 is severe imbalance.\\
		\vspace{0.5em}
		\textbf{Wrong:}
		\begin{itemize}
			\item A is balanced.
			\item C is a regression target property.
			\item D is missingness, not imbalance.
		\end{itemize}
	\end{frame}
	
	% ---------------------------------------------------------------
	\begin{frame}{Question 58}
		Which algorithm typically uses bagging?
		\begin{enumerate}[label=\Alph*.]
			\item Gradient boosted trees
			\item K-means
			\item Random forest
			\item Linear regression
		\end{enumerate}
	\end{frame}
	\begin{frame}{Answer 58}
		\textbf{Correct: C} --- Random forest uses bagging.\\
		\vspace{0.5em}
		\textbf{Wrong:}
		\begin{itemize}
			\item A uses boosting.
			\item B is clustering.
			\item D is a single model.
		\end{itemize}
	\end{frame}
	
	% ---------------------------------------------------------------
	\begin{frame}{Question 59}
		Spark ML linear regression on large data uses which approach for distributed training?
		\begin{enumerate}[label=\Alph*.]
			\item Logistic regression
			\item Spark ML cannot distribute linear regression
			\item Iterative optimization such as L-BFGS
			\item Least-squares normal equations
		\end{enumerate}
	\end{frame}
	\begin{frame}{Answer 59}
		\textbf{Correct: C} --- large data uses iterative optimizers like L-BFGS.\\
		\vspace{0.5em}
		\textbf{Wrong:}
		\begin{itemize}
			\item A is a classification algorithm.
			\item B is false.
			\item D is limited to $\le 4096$ features.
		\end{itemize}
	\end{frame}
	
	% ---------------------------------------------------------------
	\begin{frame}{Question 60}
		Why can a Spark ML decision tree differ from a scikit-learn tree with identical data and hyperparameters?
		\begin{enumerate}[label=\Alph*.]
			\item Spark ML tests every feature
			\item Spark ML automatically prunes overfit trees
			\item Spark ML tests more split candidates
			\item Spark ML tests a random sample of features
			\item Spark ML tests binned feature values as representative split candidates
		\end{enumerate}
	\end{frame}
	\begin{frame}{Answer 60}
		\textbf{Correct: E} --- Spark bins continuous features for distributed efficiency.\\
		\vspace{0.5em}
		\textbf{Wrong:}
		\begin{itemize}
			\item A is false.
			\item B pruning is user-controlled.
			\item C it tests fewer, binned candidates.
			\item D random feature sampling is for Random Forest.
		\end{itemize}
	\end{frame}
	
	% ---------------------------------------------------------------
	\begin{frame}{Question 61}
		What does Hyperopt's \texttt{fmin} do?
		\begin{enumerate}[label=\Alph*.]
			\item Minimizes an objective function over a search space
			\item Maximizes accuracy only
			\item Trains a model
			\item Logs a metric
		\end{enumerate}
	\end{frame}
	\begin{frame}{Answer 61}
		\textbf{Correct: A} --- \texttt{fmin} minimizes the objective.\\
		\vspace{0.5em}
		\textbf{Wrong:}
		\begin{itemize}
			\item B it is not limited to accuracy.
			\item C it calls a user objective.
			\item D it does not log metrics.
		\end{itemize}
	\end{frame}
	
	% ---------------------------------------------------------------
	\begin{frame}{Question 62}
		What does \texttt{SparkTrials} do?
		\begin{enumerate}[label=\Alph*.]
			\item Distributes Hyperopt trials across a Spark cluster
			\item Runs trials on the driver only
			\item Saves models to the registry
			\item Encodes categorical features
		\end{enumerate}
	\end{frame}
	\begin{frame}{Answer 62}
		\textbf{Correct: A} --- \texttt{SparkTrials} distributes trials.\\
		\vspace{0.5em}
		\textbf{Wrong:}
		\begin{itemize}
			\item B is the opposite.
			\item C is model registration.
			\item D is a preprocessing step.
		\end{itemize}
	\end{frame}
	
	% ---------------------------------------------------------------
	\begin{frame}{Question 63}
		What is the role of \texttt{tpe.suggest} in Hyperopt?
		\begin{enumerate}[label=\Alph*.]
			\item Evaluates the objective function
			\item Specifies the MLflow experiment
			\item Suggests new hyperparameter configurations based on previous results
			\item Parallelizes the tuning process
		\end{enumerate}
	\end{frame}
	\begin{frame}{Answer 63}
		\textbf{Correct: C} --- \texttt{tpe.suggest} proposes new configurations.\\
		\vspace{0.5em}
		\textbf{Wrong:}
		\begin{itemize}
			\item A is the objective function's job.
			\item B is MLflow integration.
			\item D is \texttt{SparkTrials}.
		\end{itemize}
	\end{frame}
	
	% ---------------------------------------------------------------
	\begin{frame}{Question 64}
		A team has 8 workers and 5 hyperparameters. Which approach best leverages the cluster?
		\begin{enumerate}[label=\Alph*.]
			\item Grid search on the driver
			\item Sequential manual tuning
			\item Hyperopt with \texttt{SparkTrials} and TPE
			\item Random search on the driver
		\end{enumerate}
	\end{frame}
	\begin{frame}{Answer 64}
		\textbf{Correct: C} --- Hyperopt with \texttt{SparkTrials} and TPE leverages the cluster.\\
		\vspace{0.5em}
		\textbf{Wrong:}
		\begin{itemize}
			\item A--B--D are inefficient or single-node.
		\end{itemize}
	\end{frame}
	
	% ---------------------------------------------------------------
	\begin{frame}{Question 65}
		With 3-fold CV and a grid of 3 $\times$ 2 hyperparameter values, how many models are trained in total?
		\begin{enumerate}[label=\Alph*.]
			\item 6
			\item 9
			\item 18
			\item 36
		\end{enumerate}
	\end{frame}
	\begin{frame}{Answer 65}
		\textbf{Correct: C} --- $3 \times 2 \times 3 = 18$.\\
		\vspace{0.5em}
		\textbf{Wrong:}
		\begin{itemize}
			\item A is just the combinations.
			\item B is wrong arithmetic.
			\item D doubles the folds incorrectly.
		\end{itemize}
	\end{frame}
	
	% ---------------------------------------------------------------
	\begin{frame}{Question 66}
		With 3-fold CV and a grid of 3 $\times$ 2 values, how many models can be trained in parallel?
		\begin{enumerate}[label=\Alph*.]
			\item 3
			\item 5
			\item 6
			\item 18
		\end{enumerate}
	\end{frame}
	\begin{frame}{Answer 66}
		\textbf{Correct: C} --- six combinations can run in parallel.\\
		\vspace{0.5em}
		\textbf{Wrong:}
		\begin{itemize}
			\item A is the number of folds.
			\item B is arbitrary.
			\item D is the total, not parallel.
		\end{itemize}
	\end{frame}
	
	% ---------------------------------------------------------------
	\begin{frame}{Question 67}
		What is a primary advantage of Hyperopt over manual tuning or grid search?
		\begin{enumerate}[label=\Alph*.]
			\item It efficiently explores the search space using algorithms like TPE
			\item Manual tuning is faster
			\item Grid search is more reliable
			\item Hyperopt needs more evaluations
		\end{enumerate}
	\end{frame}
	\begin{frame}{Answer 67}
		\textbf{Correct: A} --- Hyperopt explores more efficiently with TPE.\\
		\vspace{0.5em}
		\textbf{Wrong:}
		\begin{itemize}
			\item B is false.
			\item C is false.
			\item D is false.
		\end{itemize}
	\end{frame}
	
	% ---------------------------------------------------------------
	\begin{frame}{Question 68}
		When is \texttt{StringIndexer} used in a pipeline?
		\begin{enumerate}[label=\Alph*.]
			\item To convert textual or categorical values to numeric indices
			\item To impute missing values
			\item To scale numeric features
			\item To assemble vectors
		\end{enumerate}
	\end{frame}
	\begin{frame}{Answer 68}
		\textbf{Correct: A} --- StringIndexer converts categories to indices.\\
		\vspace{0.5em}
		\textbf{Wrong:}
		\begin{itemize}
			\item B is Imputer.
			\item C is StandardScaler.
			\item D is VectorAssembler.
		\end{itemize}
	\end{frame}
	
	% ---------------------------------------------------------------
	\begin{frame}{Question 69}
		Why might a train-validation split be more efficient than k-fold CV?
		\begin{enumerate}[label=\Alph*.]
			\item It requires testing fewer models, saving compute
			\item It reduces hyperparameter values to test
			\item It guarantees reproducibility
			\item It ensures bias-free training
		\end{enumerate}
	\end{frame}
	\begin{frame}{Answer 69}
		\textbf{Correct: A} --- a split trains fewer models than k-fold CV.\\
		\vspace{0.5em}
		\textbf{Wrong:}
		\begin{itemize}
			\item B is false.
			\item C is not guaranteed.
			\item D is false.
		\end{itemize}
	\end{frame}
	
	% ---------------------------------------------------------------
	\begin{frame}{Question 70}
		When using a pipeline and a cross-validator, where should they be placed?
		\begin{enumerate}[label=\Alph*.]
			\item Always put the cross-validator inside the pipeline
			\item If the pipeline includes estimators, put the entire pipeline inside the cross-validator
			\item If concerned about leakage, put the pipeline inside the cross-validator
			\item Always put the pipeline inside the cross-validator
		\end{enumerate}
	\end{frame}
	\begin{frame}{Answer 70}
		\textbf{Correct: B} --- put the pipeline inside the CV when it contains estimators.\\
		\vspace{0.5em}
		\textbf{Wrong:}
		\begin{itemize}
			\item A causes leakage.
			\item C is only the leakage special case.
			\item D is too absolute; estimators dictate placement.
		\end{itemize}
	\end{frame}
	
	% ---------------------------------------------------------------
	\begin{frame}{Question 71}
		A health organization wants to maximize the number of positive infection cases identified. Which metric should be used?
		\begin{enumerate}[label=\Alph*.]
			\item RMSE
			\item Precision
			\item Area under the residual operating curve
			\item Accuracy
			\item Recall
		\end{enumerate}
	\end{frame}
	\begin{frame}{Answer 71}
		\textbf{Correct: E} --- recall maximizes identification of positives.\\
		\vspace{0.5em}
		\textbf{Wrong:}
		\begin{itemize}
			\item A is regression.
			\item B focuses on false positives.
			\item C is not a real metric.
			\item D is misleading under imbalance.
		\end{itemize}
	\end{frame}
	
	% ---------------------------------------------------------------
	\begin{frame}{Question 72}
		When should F1 be used over accuracy?
		\begin{enumerate}[label=\Alph*.]
			\item When classes are balanced
			\item When there are more than two classes
			\item When there is significant imbalance and avoiding false negatives is a priority
			\item When true positives and true negatives are equally important
		\end{enumerate}
	\end{frame}
	\begin{frame}{Answer 72}
		\textbf{Correct: C} --- F1 is for imbalance with false-negative priority.\\
		\vspace{0.5em}
		\textbf{Wrong:}
		\begin{itemize}
			\item A accuracy is fine when balanced.
			\item B F1 extends to multiclass but is not uniquely for it.
			\item D accuracy captures both TPs and TNs.
		\end{itemize}
	\end{frame}
	
	% ---------------------------------------------------------------
	\begin{frame}{Question 73}
		A fraud dataset has a 1\% positive class and a model reports 99\% accuracy. What should you do?
		\begin{enumerate}[label=\Alph*.]
			\item Deploy it
			\item Evaluate precision, recall, F1, and PR-AUC
			\item Increase epochs
			\item Add features
		\end{enumerate}
	\end{frame}
	\begin{frame}{Answer 73}
		\textbf{Correct: B} --- evaluate imbalance-aware metrics.\\
		\vspace{0.5em}
		\textbf{Wrong:}
		\begin{itemize}
			\item A accuracy is misleading.
			\item C does not fix the metric.
			\item D is premature.
		\end{itemize}
	\end{frame}
	
	% ---------------------------------------------------------------
	\begin{frame}{Question 74}
		What does $R^2 = 0.85$ mean?
		\begin{enumerate}[label=\Alph*.]
			\item 85\% of predictions are correct
			\item The model explains 85\% of the variance in the target
			\item The model has 85\% precision
			\item The RMSE is 0.85
		\end{enumerate}
	\end{frame}
	\begin{frame}{Answer 74}
		\textbf{Correct: B} --- $R^2$ is the fraction of variance explained.\\
		\vspace{0.5em}
		\textbf{Wrong:}
		\begin{itemize}
			\item A is classification accuracy.
			\item C is precision.
			\item D is RMSE, a different metric.
		\end{itemize}
	\end{frame}
	
	% ---------------------------------------------------------------
	\begin{frame}{Question 75}
		Why might RMSE be higher than MAE?
		\begin{enumerate}[label=\Alph*.]
			\item RMSE penalizes large errors more heavily
			\item MAE squares errors
			\item RMSE ignores outliers
			\item MAE is in different units
		\end{enumerate}
	\end{frame}
	\begin{frame}{Answer 75}
		\textbf{Correct: A} --- RMSE penalizes large errors more.\\
		\vspace{0.5em}
		\textbf{Wrong:}
		\begin{itemize}
			\item B MAE uses absolute values.
			\item C RMSE is more outlier-sensitive.
			\item D both are in target units.
		\end{itemize}
	\end{frame}
	
	% ---------------------------------------------------------------
	\begin{frame}{Question 76}
		Training accuracy is 99.5\% and validation accuracy is 72\%. What is the diagnosis and fix?
		\begin{enumerate}[label=\Alph*.]
			\item Underfitting; increase depth
			\item Overfitting; reduce complexity or add regularization
			\item Good model; deploy
			\item Need more test data
		\end{enumerate}
	\end{frame}
	\begin{frame}{Answer 76}
		\textbf{Correct: B} --- a large train-validation gap means overfitting.\\
		\vspace{0.5em}
		\textbf{Wrong:}
		\begin{itemize}
			\item A underfitting shows poor train accuracy.
			\item C the gap shows poor generalization.
			\item D more test data does not fix overfitting.
		\end{itemize}
	\end{frame}
	
	% ---------------------------------------------------------------
	\begin{frame}{Question 77}
		What is the effect of increasing model complexity on bias and variance?
		\begin{enumerate}[label=\Alph*.]
			\item Bias decreases, variance increases
			\item Bias increases, variance decreases
			\item Both decrease
			\item Both increase
		\end{enumerate}
	\end{frame}
	\begin{frame}{Answer 77}
		\textbf{Correct: A} --- complexity lowers bias and raises variance.\\
		\vspace{0.5em}
		\textbf{Wrong:}
		\begin{itemize}
			\item B is the reverse.
			\item C is false.
			\item D is false.
		\end{itemize}
	\end{frame}
	
	% ---------------------------------------------------------------
	\begin{frame}{Question 78}
		A log(price) model is evaluated against actual price values and its RMSE is much larger than a price model. Which explanation is invalid?
		\begin{enumerate}[label=\Alph*.]
			\item The second model is much more accurate
			\item They failed to exponentiate predictions before computing RMSE
			\item They failed to take the log of predictions in the first model
			\item The first model is much more accurate
			\item RMSE is invalid for regression
		\end{enumerate}
	\end{frame}
	\begin{frame}{Answer 78}
		\textbf{Correct: E} --- RMSE is valid for regression.\\
		\vspace{0.5em}
		\textbf{Wrong:}
		\begin{itemize}
			\item A--D are possible explanations.
		\end{itemize}
	\end{frame}
	
	% ---------------------------------------------------------------
	\begin{frame}{Question 79}
		What is the 3-fold CV RMSE given fold RMSEs 10.0, 12.0, and 17.0?
		\begin{enumerate}[label=\Alph*.]
			\item 13.0
			\item 17.0
			\item 12.0
			\item 39.0
		\end{enumerate}
	\end{frame}
	\begin{frame}{Answer 79}
		\textbf{Correct: A} --- $(10+12+17)/3 = 13.0$.\\
		\vspace{0.5em}
		\textbf{Wrong:}
		\begin{itemize}
			\item B is the max fold.
			\item C is one fold.
			\item D is the sum.
		\end{itemize}
	\end{frame}
	
	% ---------------------------------------------------------------
	\begin{frame}{Question 80}
		What happens if the model is placed as a stage inside a pipeline and the cross-validator is the final stage?
		\begin{enumerate}[label=\Alph*.]
			\item Training takes longer
			\item Feature engineering is computed using validation data, causing leakage
			\item CV is no longer reproducible
			\item The model is refit once more per fold
		\end{enumerate}
	\end{frame}
	\begin{frame}{Answer 80}
		\textbf{Correct: B} --- feature engineering on validation data leaks.\\
		\vspace{0.5em}
		\textbf{Wrong:}
		\begin{itemize}
			\item A is not the main consequence.
			\item C CV remains reproducible.
			\item D is not the issue.
		\end{itemize}
	\end{frame}
	
	% ---------------------------------------------------------------
	\begin{frame}{Question 81}
		What does \texttt{CrossValidator} do in Spark ML?
		\begin{enumerate}[label=\Alph*.]
			\item Performs k-fold CV with a \texttt{ParamGridBuilder}
			\item Trains one model
			\item Clusters data
			\item Imputes missing values
		\end{enumerate}
	\end{frame}
	\begin{frame}{Answer 81}
		\textbf{Correct: A} --- \texttt{CrossValidator} does k-fold CV with a parameter grid.\\
		\vspace{0.5em}
		\textbf{Wrong:}
		\begin{itemize}
			\item B is a single estimator.
			\item C is clustering.
			\item D is Imputer.
		\end{itemize}
	\end{frame}
	
	% ---------------------------------------------------------------
	\begin{frame}{Question 82}
		When is a train-validation split preferable to k-fold CV?
		\begin{enumerate}[label=\Alph*.]
			\item When high interpretability is required
			\item When compute is limited, especially for large data or complex models
			\item When the algorithm is highly parallelizable
			\item When compute is unlimited
		\end{enumerate}
	\end{frame}
	\begin{frame}{Answer 82}
		\textbf{Correct: B} --- a split is cheaper than k-fold CV.\\
		\vspace{0.5em}
		\textbf{Wrong:}
		\begin{itemize}
			\item A is unrelated.
			\item C favors CV.
			\item D favors CV.
		\end{itemize}
	\end{frame}
	
	% ---------------------------------------------------------------
	\begin{frame}{Question 83}
		Which technique mitigates overfitting?
		\begin{enumerate}[label=\Alph*.]
			\item Feature engineering and dimensionality reduction
			\item Increasing model complexity
			\item Using only one technique such as data augmentation
			\item All of the above
		\end{enumerate}
	\end{frame}
	\begin{frame}{Answer 83}
		\textbf{Correct: A} --- feature engineering and dimensionality reduction mitigate overfitting.\\
		\vspace{0.5em}
		\textbf{Wrong:}
		\begin{itemize}
			\item B worsens overfitting.
			\item C is insufficient alone.
			\item D is false.
		\end{itemize}
	\end{frame}
	
	% ---------------------------------------------------------------
	\begin{frame}{Question 84}
		What is k-Nearest Neighbors' key hyperparameter tradeoff?
		\begin{enumerate}[label=\Alph*.]
			\item Small $k$ gives low bias and high variance; large $k$ gives high bias and low variance
			\item Small $k$ gives high bias; large $k$ gives high variance
			\item $k$ has no effect
			\item $k$ only affects training time
		\end{enumerate}
	\end{frame}
	\begin{frame}{Answer 84}
		\textbf{Correct: A} --- small $k$ means low bias and high variance.\\
		\vspace{0.5em}
		\textbf{Wrong:}
		\begin{itemize}
			\item B is reversed.
			\item C is false.
			\item D is false.
		\end{itemize}
	\end{frame}
	
	% ---------------------------------------------------------------
	\begin{frame}{Question 85}
		Which kernel is commonly used for non-linear SVM classification?
		\begin{enumerate}[label=\Alph*.]
			\item Linear
			\item RBF
			\item Identity
			\item Zero
		\end{enumerate}
	\end{frame}
	\begin{frame}{Answer 85}
		\textbf{Correct: B} --- RBF is common for non-linear SVM.\\
		\vspace{0.5em}
		\textbf{Wrong:}
		\begin{itemize}
			\item A is linear.
			\item C is not a kernel.
			\item D is not a kernel.
		\end{itemize}
	\end{frame}
	
	% ---------------------------------------------------------------
	\begin{frame}{Question 86}
		What is gradient boosting?
		\begin{enumerate}[label=\Alph*.]
			\item Trees trained in parallel on bootstrap samples
			\item Trees trained sequentially, each correcting previous errors
			\item A clustering method
			\item A dimensionality reduction method
		\end{enumerate}
	\end{frame}
	\begin{frame}{Answer 86}
		\textbf{Correct: B} --- boosting trains sequentially.\\
		\vspace{0.5em}
		\textbf{Wrong:}
		\begin{itemize}
			\item A is bagging.
			\item C is clustering.
			\item D is PCA.
		\end{itemize}
	\end{frame}
	
	% ---------------------------------------------------------------
	\begin{frame}{Question 87}
		What is the role of \texttt{max\_depth} in tree models?
		\begin{enumerate}[label=\Alph*.]
			\item Limits tree depth to control complexity and overfitting
			\item Sets the number of trees
			\item Sets the learning rate
			\item Sets the number of features
		\end{enumerate}
	\end{frame}
	\begin{frame}{Answer 87}
		\textbf{Correct: A} --- \texttt{max\_depth} limits depth.\\
		\vspace{0.5em}
		\textbf{Wrong:}
		\begin{itemize}
			\item B is \texttt{num\_trees}.
			\item C is \texttt{learning\_rate}.
			\item D is \texttt{max\_features}.
		\end{itemize}
	\end{frame}
	
	% ---------------------------------------------------------------
	\begin{frame}{Question 88}
		Which metric is most appropriate for a regression model with outliers?
		\begin{enumerate}[label=\Alph*.]
			\item RMSE
			\item MAE
			\item Accuracy
			\item F1
		\end{enumerate}
	\end{frame}
	\begin{frame}{Answer 88}
		\textbf{Correct: B} --- MAE is less sensitive to outliers than RMSE.\\
		\vspace{0.5em}
		\textbf{Wrong:}
		\begin{itemize}
			\item A RMSE is more outlier-sensitive.
			\item C is classification.
			\item D is classification.
		\end{itemize}
	\end{frame}
	
	% ---------------------------------------------------------------
	\begin{frame}{Question 89}
		What are the three model serving approaches tested on the exam?
		\begin{enumerate}[label=\Alph*.]
			\item Batch, real-time, and streaming
			\item Batch, offline, and online
			\item Training, validation, and testing
			\item SQL, Python, and Scala
		\end{enumerate}
	\end{frame}
	\begin{frame}{Answer 89}
		\textbf{Correct: A} --- batch, real-time, and streaming.\\
		\vspace{0.5em}
		\textbf{Wrong:}
		\begin{itemize}
			\item B uses non-exam terms.
			\item C is the ML lifecycle.
			\item D is programming languages.
		\end{itemize}
	\end{frame}
	
	% ---------------------------------------------------------------
	\begin{frame}{Question 90}
		A scikit-learn model must score 100 million Delta records efficiently. What is the best approach?
		\begin{enumerate}[label=\Alph*.]
			\item Load into pandas on the driver and loop
			\item Use \texttt{mlflow.pyfunc.spark\_udf()} to distribute inference
			\item Export to CSV and score externally
			\item Create a real-time endpoint and call per record
		\end{enumerate}
	\end{frame}
	\begin{frame}{Answer 90}
		\textbf{Correct: B} --- \texttt{spark\_udf} distributes inference.\\
		\vspace{0.5em}
		\textbf{Wrong:}
		\begin{itemize}
			\item A does not scale.
			\item C is manual.
			\item D is high-latency per record.
		\end{itemize}
	\end{frame}
	
	% ---------------------------------------------------------------
	\begin{frame}{Question 91}
		How is streaming inference performed in Databricks?
		\begin{enumerate}[label=\Alph*.]
			\item Using Delta Live Tables
			\item Using a SQL Warehouse
			\item Using the driver only
			\item Using a pandas UDF on a single node
		\end{enumerate}
	\end{frame}
	\begin{frame}{Answer 91}
		\textbf{Correct: A} --- streaming inference uses Delta Live Tables.\\
		\vspace{0.5em}
		\textbf{Wrong:}
		\begin{itemize}
			\item B is SQL only.
			\item C does not scale.
			\item D is not streaming inference.
		\end{itemize}
	\end{frame}
	
	% ---------------------------------------------------------------
	\begin{frame}{Question 92}
		How do you split traffic between real-time serving endpoints?
		\begin{enumerate}[label=\Alph*.]
			\item Configure traffic splitting on the serving endpoint
			\item Duplicate the model
			\item Use a SQL query
			\item Rename the endpoint
		\end{enumerate}
	\end{frame}
	\begin{frame}{Answer 92}
		\textbf{Correct: A} --- traffic splitting is configured on the endpoint.\\
		\vspace{0.5em}
		\textbf{Wrong:}
		\begin{itemize}
			\item B does not split traffic.
			\item C is not routing.
			\item D does not split traffic.
		\end{itemize}
	\end{frame}
	
	% ---------------------------------------------------------------
	\begin{frame}{Question 93}
		What is the primary purpose of a real-time serving endpoint?
		\begin{enumerate}[label=\Alph*.]
			\item Batch scoring
			\item Low-latency predictions accessible via REST API
			\item Training models
			\item Data storage
		\end{enumerate}
	\end{frame}
	\begin{frame}{Answer 93}
		\textbf{Correct: B} --- real-time endpoints serve low-latency REST predictions.\\
		\vspace{0.5em}
		\textbf{Wrong:}
		\begin{itemize}
			\item A is batch inference.
			\item C is training.
			\item D is storage.
		\end{itemize}
	\end{frame}
	
	% ---------------------------------------------------------------
	\begin{frame}{Question 94}
		A team replaces one model with three new models of nearly equal latency. When is the new solution less efficient during inference?
		\begin{enumerate}[label=\Alph*.]
			\item When if-else logic picks a model per prediction
			\item When average model latency exceeds the original model's size
			\item When fewer features are used
			\item When every record must be scored by all three models
		\end{enumerate}
	\end{frame}
	\begin{frame}{Answer 94}
		\textbf{Correct: D} --- scoring every record with all three multiplies cost.\\
		\vspace{0.5em}
		\textbf{Wrong:}
		\begin{itemize}
			\item A only one model runs per record.
			\item B is about model size, not inference count.
			\item C fewer features reduce cost.
		\end{itemize}
	\end{frame}
	
	% ---------------------------------------------------------------
	\begin{frame}{Question 95}
		What is a batch inference job in Databricks?
		\begin{enumerate}[label=\Alph*.]
			\item Scoring a large dataset offline and writing predictions to a table
			\item A real-time REST endpoint
			\item A training job
			\item A cluster configuration
		\end{enumerate}
	\end{frame}
	\begin{frame}{Answer 95}
		\textbf{Correct: A} --- batch inference scores offline and writes predictions.\\
		\vspace{0.5em}
		\textbf{Wrong:}
		\begin{itemize}
			\item B is real-time serving.
			\item C is training.
			\item D is compute setup.
		\end{itemize}
	\end{frame}
	
	% ---------------------------------------------------------------
	\begin{frame}{Question 96}
		Which compute type is best for Spark SQL data prep plus Spark ML training in a notebook?
		\begin{enumerate}[label=\Alph*.]
			\item Single Node cluster
			\item Standard (all-purpose) cluster
			\item SQL Warehouse
			\item None of these
		\end{enumerate}
	\end{frame}
	\begin{frame}{Answer 96}
		\textbf{Correct: B} --- a Standard cluster supports SQL and Spark ML.\\
		\vspace{0.5em}
		\textbf{Wrong:}
		\begin{itemize}
			\item A is single-node.
			\item C is SQL-only.
			\item D is false.
		\end{itemize}
	\end{frame}
	
	% ---------------------------------------------------------------
	\begin{frame}{Question 97}
		Why is model monitoring important in production?
		\begin{enumerate}[label=\Alph*.]
			\item Data distributions and relationships can drift over time
			\item Once deployed, models never degrade
			\item It replaces evaluation
			\item It improves training speed
		\end{enumerate}
	\end{frame}
	\begin{frame}{Answer 97}
		\textbf{Correct: A} --- drift degrades deployed models.\\
		\vspace{0.5em}
		\textbf{Wrong:}
		\begin{itemize}
			\item B is false.
			\item C evaluation is separate.
			\item D monitoring is runtime, not training.
		\end{itemize}
	\end{frame}
	
	% ---------------------------------------------------------------
	\begin{frame}{Question 98}
		What is data drift versus concept drift?
		\begin{enumerate}[label=\Alph*.]
			\item Data drift changes input distributions; concept drift changes the feature-target relationship
			\item Both refer to the same thing
			\item Data drift changes labels; concept drift changes features
			\item Neither affects models
		\end{enumerate}
	\end{frame}
	\begin{frame}{Answer 98}
		\textbf{Correct: A} --- data drift changes inputs; concept drift changes the relationship.\\
		\vspace{0.5em}
		\textbf{Wrong:}
		\begin{itemize}
			\item B is false.
			\item C reverses the definitions.
			\item D is false.
		\end{itemize}
	\end{frame}
	
	% ---------------------------------------------------------------
	\begin{frame}{Question 99}
		Which is a fairness concern when deploying ML?
		\begin{enumerate}[label=\Alph*.]
			\item Biased training data leading to discriminatory outcomes
			\item Faster inference
			\item Lower storage cost
			\item Better dashboards
		\end{enumerate}
	\end{frame}
	\begin{frame}{Answer 99}
		\textbf{Correct: A} --- biased data can cause discriminatory outcomes.\\
		\vspace{0.5em}
		\textbf{Wrong:}
		\begin{itemize}
			\item B, C, D are performance or cost concerns, not fairness.
		\end{itemize}
	\end{frame}
	
	% ---------------------------------------------------------------
	\begin{frame}{Question 100}
		Which practice best guards against leakage in a production pipeline?
		\begin{enumerate}[label=\Alph*.]
			\item Fit all transforms on the entire dataset
			\item Encapsulate preprocessing and model in a Pipeline fit only on training data
			\item Use the test set for feature selection
			\item Scale before splitting
		\end{enumerate}
	\end{frame}
	\begin{frame}{Answer 100}
		\textbf{Correct: B} --- a Pipeline fit on training data prevents leakage.\\
		\vspace{0.5em}
		\textbf{Wrong:}
		\begin{itemize}
			\item A leaks test information.
			\item C uses test data for selection.
			\item D scales before splitting, leaking.
		\end{itemize}
	\end{frame}
	
\end{document}